\documentclass[a4paper,12pt]{article}
\usepackage[T2A]{fontenc}
\usepackage[utf8]{inputenc}
\usepackage[russian]{babel}
\usepackage{amsmath,amssymb,mathtools}
\usepackage{icomma}
\usepackage[left=23mm,right=23mm,top=23mm,bottom=23mm,headheight=24pt,headsep=8mm]{geometry}
\usepackage{microtype}
\usepackage{tikz}
\usetikzlibrary{arrows.meta,positioning,shapes.geometric,calc}
\usepackage{tabularx,booktabs}
\usepackage{enumitem}
\usepackage{fancyhdr}
\usepackage{xcolor}
\usepackage{setspace}
\usepackage{parskip}
\usepackage{needspace}

\definecolor{accent}{HTML}{008080}
\definecolor{darkaccent}{HTML}{005757}
\definecolor{textgray}{HTML}{303637}
\color{textgray}
\setstretch{1.18}
\setlength{\parindent}{0pt}
\setlength{\parskip}{0.45em}
\setlist[itemize]{leftmargin=1.55em,itemsep=0.2em,topsep=0.3em}
\setlist[enumerate]{leftmargin=1.7em,itemsep=0.25em,topsep=0.35em}
\pagestyle{fancy}
\fancyhf{}
\fancyhead[L]{\footnotesize\color{textgray} Метод интервалов и преобразование неравенств}
\fancyhead[R]{\footnotesize\color{textgray} Ярослав Верещагин}
\fancyfoot[C]{\footnotesize\thepage}
\renewcommand{\headrulewidth}{0.45pt}
\renewcommand{\headrule}{\hbox to\headwidth{\color{accent}\leaders\hrule height \headrulewidth\hfill}}
\newcommand{\sectiontitle}[1]{\needspace{4\baselineskip}\vspace{0.55em}{\large\bfseries\color{darkaccent}#1}\par\vspace{0.15em}\noindent{\color{accent}\rule{\linewidth}{0.45pt}}\par\vspace{0.2em}}
\newcommand{\minititle}[1]{\needspace{3\baselineskip}{\bfseries\color{darkaccent}#1}\par}
\newcommand{\ans}[1]{\par\textbf{Ответ:} $#1$.}

\begin{document}
\begin{titlepage}
\thispagestyle{empty}
\centering
\vspace*{24mm}
{\large\color{darkaccent} МАТЕМАТИКА\par}
\vspace{24mm}
{\LARGE Чек-лист\par}
\vspace{12mm}
{\LARGE\color{accent} Метод интервалов\par}
\vspace{2mm}
{\Large\color{accent} и преобразование неравенств\par}
\vspace{21mm}
{\normalsize Подготовка к ОГЭ по математике\par}
\vfill
{\large Ярослав Верещагин\par}
\vspace{12mm}
{\normalsize 29 октября 2026 г.\par}
\vspace{19mm}
{\small Лёвин Артём Александрович\par}
{\small эксклюзивно для Ярослава Верещагина\par}
\vspace{17mm}
\begin{tikzpicture}
\draw[accent,line width=0.95pt] (-5.1,0) .. controls (-3.7,0.65) and (-1.9,-0.55) .. (0,0)
                                  .. controls (1.7,0.53) and (3.65,-0.58) .. (5.1,0);
\end{tikzpicture}
\vspace*{6mm}
\end{titlepage}
\setcounter{page}{1}

\sectiontitle{1. Нормализация --- первый шаг}
Метод интервалов удобно применять, когда справа стоит нуль:
\[
F(x)>0,\qquad F(x)<0,\qquad F(x)\ge0,\qquad F(x)\le0.
\]
\textbf{Нормализовать неравенство} --- перенести все слагаемые в одну сторону и упростить полученное выражение. При переносе через знак неравенства меняется знак \emph{слагаемого}; сам знак неравенства сохраняется.

\minititle{Пример: сначала раскрываем скобки}
\[
\begin{aligned}
(x-3)(2x+3)&<-7,\\
(x-3)(2x+3)+7&<0,\\
2x^2+3x-6x-9+7&<0,\\
2x^2-3x-2&<0.
\end{aligned}
\]
Нули квадратного трёхчлена находим из \emph{уравнения}
\[
2x^2-3x-2=0,\quad D=(-3)^2-4\cdot2\cdot(-2)=25,
\quad x_1=-\frac12,\quad x_2=2.
\]
Коэффициент при $x^2$ положительный. Между двумя простыми корнями выражение отрицательно. Границы выколоты: знак исходного неравенства строгий.
\ans{\left(-\frac12;2\right)}

\minititle{Не перепутайте знаки}
\begin{itemize}
\item В строках преобразования остаётся знак $<$, $>$, $\le$ или $\ge$.
\item Знак $=$ пишется при поиске нулей числителя, знаменателя или многочлена.
\item Пробную точку подставляют в \textbf{нормализованное выражение}.
\end{itemize}

\sectiontitle{2. Общий множитель: сохраняем скобки}
Когда одно и то же выражение повторяется, вынесите его за скобки:
\[
\begin{aligned}
(x-7)^2&<\sqrt{11}\,(x-7),\\
(x-7)^2-\sqrt{11}\,(x-7)&<0,\\
(x-7)\bigl((x-7)-\sqrt{11}\bigr)&<0.
\end{aligned}
\]
Нули множителей: $x=7$ и $x=7+\sqrt{11}$. Оба простые, поэтому между ними произведение отрицательно.
\ans{(7;7+\sqrt{11})}
\textbf{Проверка:} при вынесении множителя $x-7$ оставшаяся скобка записывается полностью. Терять внутренние скобки нельзя.

\sectiontitle{3. Знаменатель и область допустимых значений}
Для дроби $\dfrac{P(x)}{Q(x)}$ необходимо $Q(x)\ne0$. Даже при нестрогом знаке неравенства \textbf{нули знаменателя всегда исключаются}.

\minititle{Пример: постоянный числитель}
\[
\frac{-10}{(x-3)^2-5}\ge0.
\]
Числитель $-10$ никогда не равен нулю. Знаменатель равен нулю при
\[
(x-3)^2=5\quad\Longrightarrow\quad
x=3-\sqrt5\quad\text{или}\quad x=3+\sqrt5.
\]
Дробь положительна там, где знаменатель отрицателен:
\[
(x-3)^2-5<0\quad\Longleftrightarrow\quad
3-\sqrt5<x<3+\sqrt5.
\]
Несмотря на знак $\ge$, обе границы выколоты, так как обращают знаменатель в нуль.
\ans{(3-\sqrt5;3+\sqrt5)}

\minititle{Как расположить корни с радикалами}
Число $3-\sqrt5$ меньше числа $3+\sqrt5$, поэтому первое располагаем слева. Для выбора пробной точки полезно знать, что $2<\sqrt5<3$: значит, $3-\sqrt5>0$, а число $0$ лежит левее обоих корней.

\sectiontitle{4. Если в неравенстве несколько дробей}
Приведите дроби к общему знаменателю и объедините их в одну дробь. Ограничения исходных знаменателей запишите \textbf{до сокращения}.

\minititle{Пример: общий знаменатель}
\[
\begin{aligned}
\frac1x-\frac1{x-4}&>0,\qquad x\ne0,\quad x\ne4,\\[2pt]
\frac{x-4-x}{x(x-4)}&>0,\\[2pt]
\frac{-4}{x(x-4)}&>0.
\end{aligned}
\]
Числитель отрицателен при любом $x$. Поэтому дробь положительна, когда $x(x-4)<0$, то есть между $0$ и $4$.
\ans{(0;4)}

\minititle{Правило «ВУПС»: скобки берегут знаки}
\textbf{Вычитаемое, умножаемое, подставляемое --- в скобки}, если внутри есть сумма или разность. Например,
\[
\frac{x-1}{x}-\frac{x-3}{x-2}
=\frac{(x-1)(x-2)-x(x-3)}{x(x-2)}.
\]
При раскрытии скобок минус перед $x(x-3)$ действует \emph{на все слагаемые}: $-x(x-3)=-x^2+3x$. Пропущенный минус или скобка меняет ответ.

\sectiontitle{5. Критические точки и знаки}
\begin{enumerate}
\item Найдите действительные нули числителя и знаменателя и расположите их по возрастанию.
\item Отметьте нули знаменателя выколотыми точками; нули числителя включайте только при нестрогом знаке и допустимом знаменателе.
\item Вычислите знак дроби на одном из промежутков, подставив удобную пробную точку.
\item Через нуль \textbf{нечётной} кратности знак меняется, через нуль \textbf{чётной} кратности сохраняется.
\item Выберите промежутки со знаком, который требуется в неравенстве.
\end{enumerate}
\textbf{Важно:} правило чередования знаков применяют с учётом кратностей. Множитель $(x-a)^2$ не меняет знак при переходе через $a$.

\sectiontitle{6. Замена повторяющегося выражения}
Если знаменатели отличаются лишь на постоянное число, полезна замена. Она помогает избежать раскрытия скобок и появления четвёртой степени.

\minititle{Пример: сохраняем структуру квадратных выражений}
\[
\frac1{x^2+4x+3}-\frac1{x^2+4x+4}>\frac12.
\]
Обозначим $t=x^2+4x+3$. Тогда $x^2+4x+4=t+1$ и
\[
\frac1t-\frac1{t+1}>\frac12,\qquad t\ne0,\quad t\ne-1.
\]
Приведение к общему знаменателю выполняем \emph{в новой переменной}:
\[
\begin{aligned}
\frac1t-\frac1{t+1}-\frac12&>0,\\[3pt]
\frac{2(t+1)-2t-t(t+1)}{2t(t+1)}&>0,\\[3pt]
\frac{(1-t)(t+2)}{2t(t+1)}&>0.
\end{aligned}
\]
\textbf{Не забывайте ограничения замены.} Поскольку
\[
t=(x+2)^2-1\ge-1,
\]
при допустимом $t$ множители $t+1$ и $t+2$ положительны. Неравенство выполняется при $0<t<1$.

\minititle{Как вернуться к переменной $x$}
\[
\begin{aligned}
0<(x+2)^2-1<1
&\quad\Longleftrightarrow\quad 1<(x+2)^2<2,\\
&\quad\Longleftrightarrow\quad
\begin{cases}
-\sqrt2<x+2<-1,\\[-1pt]
\text{или}\\[-1pt]
1<x+2<\sqrt2.
\end{cases}
\end{aligned}
\]
\ans{(-2-\sqrt2;-3)\cup(-1;-2+\sqrt2)}

\vspace{1mm}
\textbf{Краткий ориентир:} сначала ограничения и решение для $t$, затем возврат к $x$ и проверка концов промежутков.

\clearpage
\begin{center}
{\Large\bfseries\color{darkaccent} Алгоритм: от исходного неравенства к ответу}
\end{center}
\vspace{5mm}
\begin{center}
\begin{tikzpicture}[
  >={Stealth[length=2.4mm]},
  term/.style={ellipse,draw=darkaccent,line width=.8pt,minimum width=34mm,minimum height=10mm,align=center,font=\small},
  proc/.style={rounded corners=2pt,draw=accent,line width=.8pt,text width=112mm,minimum height=12mm,inner sep=5pt,align=center,font=\small},
  decide/.style={diamond,draw=darkaccent,line width=.8pt,aspect=3.4,text width=54mm,inner sep=1pt,align=center,font=\small},
  branch/.style={rounded corners=2pt,draw=accent,line width=.8pt,text width=50mm,minimum height=16mm,inner sep=5pt,align=center,font=\small},
  arr/.style={->,line width=.8pt,draw=textgray}
]
\node[term] (start) {Начало};
\node[proc,below=9mm of start] (norm) {Перенести всё влево; справа получить $0$; сохранить знак неравенства};
\node[decide,below=11mm of norm] (sub) {Есть удобное повторяющееся выражение?};
\node[branch,below left=13mm and 1mm of sub] (yes) {Сделать замену $t$, записать её ограничения};
\node[branch,below right=13mm and 1mm of sub] (no) {Оставить переменную $x$};
\node[proc,below=17mm of $(yes.south)!0.5!(no.south)$] (frac) {Раскрыть нужные скобки или вынести общий множитель; дроби привести к общему знаменателю};
\node[proc,below=8mm of frac] (roots) {Найти нули числителя и знаменателя; нули знаменателя исключить};
\node[proc,below=8mm of roots] (signs) {Отметить точки по порядку, определить знаки, выбрать нужные промежутки};
\node[proc,below=8mm of signs] (back) {Если была замена, вернуться к $x$; проверить ОДЗ и скобки};
\node[term,below=9mm of back] (finish) {Ответ};
\draw[arr] (start)--(norm);
\draw[arr] (norm)--(sub);
\draw[arr] (sub.south west) -| node[pos=.18,above left,font=\footnotesize]{да} (yes.north);
\draw[arr] (sub.south east) -| node[pos=.18,above right,font=\footnotesize]{нет} (no.north);
\draw[arr] (yes.south) |- (frac.north);
\draw[arr] (no.south) |- (frac.north);
\draw[arr] (frac)--(roots);
\draw[arr] (roots)--(signs);
\draw[arr] (signs)--(back);
\draw[arr] (back)--(finish);
\end{tikzpicture}
\end{center}

\clearpage
\sectiontitle{Тренировочный блок: Путь А}
Решите \textbf{10 неравенств}. Знаки при преобразованиях сохраняйте, нули знаменателя исключайте. В задачах с заменой покажите переход к новой переменной и возврат к $x$.

\minititle{Базовый уровень}
\begin{enumerate}[label=\arabic*.,itemsep=.8em]
\item $\displaystyle (x-4)(x+1)<6.$
\item $\displaystyle (x-5)^2<3(x-5).$
\item $\displaystyle (x-6)^2\ge\sqrt7\,(x-6).$
\end{enumerate}

\minititle{Уровень выше среднего}
\begin{enumerate}[label=\arabic*.,start=4,itemsep=.8em]
\item $\displaystyle \frac{12}{(x-2)^2-9}\le0.$
\item $\displaystyle \frac1{x-1}-\frac1{x+3}\ge0.$
\item $\displaystyle \frac1x-\frac1{x-5}>0.$
\item $\displaystyle \frac1{x-2}-\frac{x-1}{x}\ge0.$
\end{enumerate}

\minititle{Сложный уровень}
\begin{enumerate}[label=\arabic*.,start=8,itemsep=.8em]
\item $\displaystyle \frac{x-2}{x-3}-\frac{x-1}{x+1}>0.$
\item $\displaystyle \frac1{x^2+2x}-\frac1{x^2+2x+1}>\frac12.$
\item $\displaystyle \frac1{x^2+2x-3}+\frac23\le\frac1{x^2+2x}.$
\end{enumerate}

\vspace{6mm}
\minititle{Самопроверка оформления}
\begin{itemize}
\item При переносе изменили знак слагаемого, но сохранили знак неравенства?
\item Умножаемые и вычитаемые выражения с суммой или разностью заключены в скобки?
\item Исключили все нули знаменателей, включая те, что могли сократиться?
\item Отличили строгие границы от включённых точек числителя?
\item После замены получили ответ именно для $x$?
\end{itemize}

\clearpage
\sectiontitle{Ответы к тренировочному блоку}
\vspace{1.5mm}
\renewcommand{\arraystretch}{1.62}
\begin{tabularx}{\linewidth}{@{}p{15mm}X@{}}
\toprule
\textbf{№} & \textbf{Ответ}\\
\midrule
1 & $(-2;5)$\\
2 & $(5;8)$\\
3 & $(-\infty;6]\cup[6+\sqrt7;+\infty)$\\
4 & $(-1;5)$\\
5 & $(-\infty;-3)\cup(1;+\infty)$\\
6 & $(0;5)$\\
7 & $(0;2-\sqrt2]\cup(2;2+\sqrt2]$\\
8 & $(-1;\frac53)\cup(3;+\infty)$\\
9 & $(-1-\sqrt2;-2)\cup(0;-1+\sqrt2)$\\
10 & $(-3;-2)\cup(0;1)$\\
\bottomrule
\end{tabularx}

\vspace{10mm}
\sectiontitle{Итоговая памятка}
\renewcommand{\arraystretch}{1.35}
\begin{tabularx}{\linewidth}{@{}p{.35\linewidth}X@{}}
\toprule
\textbf{Что видно в задаче} & \textbf{Что сделать}\\
\midrule
Справа не нуль & Перенести все слагаемые влево, сохранив знак неравенства.\\
Повторяющийся множитель & Вынести общую скобку, сохранив её структуру.\\
Несколько дробей & Привести к общему знаменателю и соблюсти правило скобок.\\
Нуль знаменателя & Исключить из ОДЗ и из ответа при любом знаке.\\
Нули на числовой оси & Разместить по возрастанию, проверить знаки и кратности.\\
Похожие квадратные знаменатели & Сделать замену всего повторяющегося выражения, затем вернуться к $x$.\\
\bottomrule
\end{tabularx}
\vfill
\begin{center}
{\footnotesize\color{darkaccent} Лёвин Артём Александрович --- эксклюзивно для Ярослава Верещагина}
\end{center}
\end{document}
